% CORRECTED VERSION
% Changes:
% - Added a stability assumption η_max ≤ 1/L and a lemma giving the lower bound
%   η - (L/2)η^2 ≥ η/2 for η ≤ 1/L (fixes Step 6).
% - Revised the descent inequality in Step 10 to keep the exact factor Lη_max^2/η_min,
%   and introduced an explicit constant‑ratio condition to recover the simpler
%   Lη_max term (fixes Step 10).
% - Removed the unused bounded‑gradient assumption and clarified the use of all
%   remaining assumptions.
% - Kept all original step annotations for easy reference.

\documentclass[11pt]{article}
\usepackage{amsmath,amssymb,amsthm}
\usepackage{algorithm,algorithmic}
\usepackage{enumitem}
\usepackage{hyperref}
\usepackage{geometry}
\geometry{margin=1in}
\setlist[enumerate]{leftmargin=*}

\newtheorem{theorem}{Theorem}
\newtheorem{lemma}{Lemma}
\newtheorem{assumption}{Assumption}

\begin{document}

\section*{Cyclic Stochastic Gradient Descent (Cyclic SGD)}

\section*{Mathematical Formulation}
Let $f:\mathbb{R}^d\rightarrow\mathbb{R}$ be the objective function and let
$\{z_t\}_{t\ge 0}$ be an i.i.d.\ data stream.
At iteration $t$ a stochastic gradient 
\[
g_t\;:=\;\nabla f(w_t;z_t)
\]
is computed.  
The learning rate follows a \emph{triangular cyclic schedule} with period $M\in\mathbb{N}$,
minimum step size $\eta_{\min}>0$ and maximum step size $\eta_{\max}>\eta_{\min}$:
\[
\eta_t \;=\;
\begin{cases}
\displaystyle 
\eta_{\min}+ \frac{2(\eta_{\max}-\eta_{\min})}{M}\,t \quad & 0\le t <\frac{M}{2},\\[6pt]
\displaystyle 
\eta_{\max}- \frac{2(\eta_{\max}-\eta_{\min})}{M}\,\Bigl(t-\frac{M}{2}\Bigr) \quad & \frac{M}{2}\le t < M,
\end{cases}
\tag{1}
\]
and the pattern repeats every $M$ steps, i.e. $\eta_{t+M}=\eta_t$.

The update rule is
\[
w_{t+1}=w_t-\eta_t\,g_t,\qquad t=0,1,2,\dots .
\tag{2}
\]

\section*{Pseudocode}
\begin{algorithm}[H]
\caption{Cyclic SGD}
\begin{algorithmic}[1]
\REQUIRE Initial point $w_0\in\mathbb{R}^d$, period $M$, $\eta_{\min},\eta_{\max}$, number of iterations $T$.
\FOR{$t=0$ \TO $T-1$}
    \STATE Compute stochastic gradient $g_t=\nabla f(w_t;z_t)$.
    \STATE Compute learning rate $\eta_t$ according to (1).
    \STATE Update $w_{t+1}=w_t-\eta_t g_t$.
\ENDFOR
\RETURN $w_T$ (or optionally the average $\bar w_T=\frac1T\sum_{t=0}^{T-1}w_t$).
\end{algorithmic}
\end{algorithm}

\section*{Dry Run Example}
Consider the scalar quadratic $f(w)=(w-3)^2$, thus $\nabla f(w)=2(w-3)$.
Choose $w_0=0$, $M=4$, $\eta_{\min}=0.05$, $\eta_{\max}=0.15$, and run three iterations.

\begin{enumerate}
\item \textbf{Iteration $t=0$:} 
\[
\eta_0=\eta_{\min}=0.05,\qquad
g_0=2(w_0-3)=2(0-3)=-6.
\]
Update:
\[
w_1=w_0-\eta_0 g_0=0-0.05(-6)=0.30.
\]
Objective:
\[
f(w_1)=(0.30-3)^2=7.29.
\]

\item \textbf{Iteration $t=1$:} 
\[
\eta_1=\eta_{\min}+ \frac{2(\eta_{\max}-\eta_{\min})}{M}\cdot1
       =0.05+ \frac{2(0.10)}{4}=0.10,
\]
\[
g_1=2(w_1-3)=2(0.30-3)=-5.40.
\]
Update:
\[
w_2=w_1-\eta_1 g_1=0.30-0.10(-5.40)=0.84.
\]
Objective:
\[
f(w_2)=(0.84-3)^2=4.66.
\]

\item \textbf{Iteration $t=2$:} 
\[
\eta_2=\eta_{\max}=0.15,
\]
\[
g_2=2(w_2-3)=2(0.84-3)=-4.32.
\]
Update:
\[
w_3=w_2-\eta_2 g_2=0.84-0.15(-4.32)=1.50.
\]
Objective:
\[
f(w_3)=(1.50-3)^2=2.25.
\]
\end{enumerate}
The example demonstrates decreasing objective values while the step size cycles.

\newpage
\section*{Assumptions}
\begin{assumption}[Smoothness]\label{as:smooth}
$f$ is $L$‑smooth: for all $x,y\in\mathbb{R}^d$,
\[
f(y)\le f(x)+\langle\nabla f(x),y-x\rangle+\frac{L}{2}\|y-x\|^2 .
\]
\end{assumption}
\begin{assumption}[Bounded Variance]\label{as:var}
For every $w$, the stochastic gradient satisfies
\[
\mathbb{E}[g_t\mid w_t]=\nabla f(w_t),\qquad
\mathbb{E}\big[\|g_t-\nabla f(w_t)\|^2\mid w_t\big]\le\sigma^2 .
\]
\end{assumption}
\begin{assumption}[Learning‑rate bounds and stability]\label{as:lr}
The cyclic schedule satisfies
\[
0<\eta_{\min}\le\eta_t\le\eta_{\max},\qquad\forall t,
\]
the period $M$ is a fixed positive integer, and the maximal step size respects the
standard stability condition
\[
\eta_{\max}\le\frac{1}{L}.
\tag{3}
\]
\end{assumption}
\begin{assumption}[Constant ratio]\label{as:ratio}
There exists a constant $c\ge1$ such that
\[
\eta_{\max}=c\,\eta_{\min}.
\tag{4}
\]
\end{assumption}
(The bounded‑gradient assumption from the original manuscript is omitted because it is never used.)

\section*{Main Convergence Theorem}
\begin{theorem}\label{thm:main}
Assume \ref{as:smooth}–\ref{as:lr} hold and let $w^*$ be a global minimizer of $f$.
Define $T$ total iterations and $\bar w_T:=\frac1T\sum_{t=0}^{T-1}w_t$.
Then
\[
\frac1T\sum_{t=0}^{T-1}\mathbb{E}\big[\|\nabla f(w_t)\|^2\big]
\;\le\;
\frac{2\big(f(w_0)-f(w^*)\big)}{\eta_{\min}T}
\;+\;
\frac{L\,\eta_{\max}^{2}}{\eta_{\min}}\;\sigma^{2}.
\tag{5}
\]
Consequently, under the constant‑ratio assumption \ref{as:ratio} (i.e. $\eta_{\max}=c\eta_{\min}$) we obtain
\[
\frac1T\sum_{t=0}^{T-1}\mathbb{E}\big[\|\nabla f(w_t)\|^2\big]
\;=\;
O\!\left(\frac{1}{\sqrt{T}}\right)
\quad\text{provided }\eta_{\max}=O\!\left(\frac{1}{\sqrt{T}}\right).
\]
\end{theorem}

\section*{Theoretical Properties}
\begin{enumerate}
\item \textbf{Stability.} Because every $\eta_t\le\eta_{\max}\le1/L$, the standard $L$‑smooth descent lemma guarantees that the iterates remain in a bounded region provided $w_0$ is bounded.
\item \textbf{Hyper‑parameter sensitivity.} The bound \eqref{5} depends on the ratio $c=\eta_{\max}/\eta_{\min}$ only through the second term $L c\eta_{\max}\sigma^{2}$ (after invoking \ref{as:ratio}). Choosing $c$ too large inflates the first term via a smaller $\eta_{\min}$, while a too small $c$ reduces the benefit of occasional larger steps.
\item \textbf{Comparison to standard SGD.} For constant step size $\eta$, the classic non‑convex SGD bound is
\[
\frac1T\sum_{t=0}^{T-1}\mathbb{E}\|\nabla f(w_t)\|^2\le
\frac{2\big(f(w_0)-f(w^*)\big)}{\eta T}+L\eta\sigma^2 .
\]
Cyclic SGD replaces the constant $\eta$ by the *average* of the cycle,
\[
\bar\eta:=\frac{1}{M}\sum_{t=0}^{M-1}\eta_t=\frac{\eta_{\min}+\eta_{\max}}{2},
\]
and the bound \eqref{5} shows that the worst‑case $\eta_{\max}$ appears only in the variance term, leading to a potentially tighter guarantee when $\eta_{\max}>\bar\eta$ but $\eta_{\min}$ is not too small.
\end{enumerate}

\section*{Discussion}
The theorem demonstrates that a cyclic learning rate does not degrade the standard $O(1/\sqrt{T})$ rate for non‑convex stochastic optimization. The key insight is that the descent contribution is governed by the *minimum* step size (the most conservative update), whereas the variance amplification is controlled by the *maximum* step size. By carefully selecting $\eta_{\min}$ and $\eta_{\max}$ (and keeping their ratio bounded), practitioners can obtain aggressive exploration (large $\eta_{\max}$) without sacrificing the asymptotic convergence guarantee.

\newpage
\section*{Preliminaries (Definitions and Lemmas)}
\begin{lemma}[Smoothness Descent Inequality]\label{lem:smooth}
If $f$ is $L$‑smooth, then for any $w$ and any stochastic gradient $g$,
\[
f(w-\eta g)\le f(w)-\eta\langle\nabla f(w),g\rangle+\frac{L\eta^{2}}{2}\|g\|^{2}.
\]
\end{lemma}
\begin{proof}
Apply the definition of $L$‑smoothness with $y=w-\eta g$ and rearrange terms. \hfill $\square$
\end{proof}

\begin{lemma}[Variance Decomposition]\label{lem:variance}
For any $w$,
\[
\mathbb{E}\big[\|g-\nabla f(w)\|^{2}\mid w\big]
=\mathbb{E}\big[\|g\|^{2}\mid w\big]-\|\nabla f(w)\|^{2}.
\]
\end{lemma}
\begin{proof}
Expand $\|g\|^{2}=\|g-\nabla f(w)+\nabla f(w)\|^{2}$ and use $\mathbb{E}[g\mid w]=\nabla f(w)$. \hfill $\square$
\end{proof}

\begin{lemma}[Lower bound for $\eta-\frac{L}{2}\eta^{2}$]\label{lem:lb}
If $0<\eta\le\frac{1}{L}$ then
\[
\eta-\frac{L}{2}\eta^{2}\;\ge\;\frac{\eta}{2}.
\]
\end{lemma}
\begin{proof}
Consider the function $h(\eta)=\eta-\frac{L}{2}\eta^{2}-\frac{\eta}{2}
=\frac{\eta}{2}\bigl(1-L\eta\bigr)$.  
Since $\eta\le 1/L$, the factor $(1-L\eta)\ge0$, hence $h(\eta)\ge0$, which yields the claimed inequality. \hfill $\square$
\end{proof}

\section*{Main Proof (Step‑by‑Step Derivation)}
We prove Theorem \ref{thm:main}. Throughout the proof, expectations are taken with respect to the randomness up to iteration $t$ (i.e., conditioning on $w_t$).

\begin{enumerate}
\item[Step 1.] Apply Lemma \ref{lem:smooth} with $w=w_t$, $g=g_t$, and $\eta=\eta_t$:
\[
f(w_{t+1})\le f(w_t)-\eta_t\langle\nabla f(w_t),g_t\rangle+\frac{L\eta_t^{2}}{2}\|g_t\|^{2}.
\tag{6}
\]

\item[Step 2.] Take conditional expectation $\mathbb{E}[\cdot\mid w_t]$ on both sides of \eqref{6}. Using unbiasedness $\mathbb{E}[g_t\mid w_t]=\nabla f(w_t)$,
\[
\mathbb{E}\big[f(w_{t+1})\mid w_t\big]\le
f(w_t)-\eta_t\|\nabla f(w_t)\|^{2}
+\frac{L\eta_t^{2}}{2}\,\mathbb{E}\big[\|g_t\|^{2}\mid w_t\big].
\tag{7}
\]

\item[Step 3.] Decompose $\mathbb{E}[\|g_t\|^{2}\mid w_t]$ using Lemma \ref{lem:variance}:
\[
\mathbb{E}\big[\|g_t\|^{2}\mid w_t\big]
= \|\nabla f(w_t)\|^{2}+\mathbb{E}\big[\|g_t-\nabla f(w_t)\|^{2}\mid w_t\big]
\le \|\nabla f(w_t)\|^{2}+\sigma^{2},
\tag{8}
\]
where the inequality follows from Assumption \ref{as:var}.

\item[Step 4.] Substitute \eqref{8} into \eqref{7}:
\[
\mathbb{E}\big[f(w_{t+1})\mid w_t\big]\le
f(w_t)-\eta_t\|\nabla f(w_t)\|^{2}
+\frac{L\eta_t^{2}}{2}\big(\|\nabla f(w_t)\|^{2}+\sigma^{2}\big).
\tag{9}
\]

\item[Step 5.] Rearrange \eqref{9} to isolate the gradient term:
\[
\big(\eta_t-\tfrac{L}{2}\eta_t^{2}\big)\,\|\nabla f(w_t)\|^{2}
\;\le\;
f(w_t)-\mathbb{E}[f(w_{t+1})\mid w_t]
+\tfrac{L}{2}\eta_t^{2}\sigma^{2}.
\tag{10}
\]

\item[Step 6.] % CORRECTED: proper lower bound using Lemma \ref{lem:lb}
Using Lemma \ref{lem:lb} and the stability condition \eqref{3},
\[
\eta_t-\frac{L}{2}\eta_t^{2}\;\ge\;\frac{\eta_t}{2}\;\ge\;\frac{\eta_{\min}}{2}.
\tag{11}
\]

\item[Step 7.] Plugging \eqref{11} into \eqref{10} yields
\[
\frac{\eta_{\min}}{2}\,\|\nabla f(w_t)\|^{2}
\;\le\;
f(w_t)-\mathbb{E}[f(w_{t+1})\mid w_t]
+\frac{L}{2}\eta_{\max}^{2}\sigma^{2}.
\tag{12}
\]

\item[Step 8.] Take total expectation over the history up to iteration $t$ and sum \eqref{12} for $t=0,\dots,T-1$:
\[
\frac{\eta_{\min}}{2}\sum_{t=0}^{T-1}\mathbb{E}\big[\|\nabla f(w_t)\|^{2}\big]
\le
\mathbb{E}[f(w_0)]-\mathbb{E}[f(w_T)]
+ \frac{L}{2}\eta_{\max}^{2}\sigma^{2}T .
\tag{13}
\]

\item[Step 9.] Since $f(w_T)\ge f(w^*)$, we obtain
\[
\frac{\eta_{\min}}{2}\sum_{t=0}^{T-1}\mathbb{E}\big[\|\nabla f(w_t)\|^{2}\big]
\le
f(w_0)-f(w^*) + \frac{L}{2}\eta_{\max}^{2}\sigma^{2}T .
\tag{14}
\]

\item[Step 10.] % ADDED: explicit division and use of constant‑ratio assumption
Divide both sides of \eqref{14} by $T$ and by $\eta_{\min}/2$:
\[
\frac1T\sum_{t=0}^{T-1}\mathbb{E}\big[\|\nabla f(w_t)\|^{2}\big]
\le
\frac{2\big(f(w_0)-f(w^*)\big)}{\eta_{\min}T}
+ \frac{L\,\eta_{\max}^{2}}{\eta_{\min}}\;\sigma^{2}.
\tag{15}
\]
This is exactly the bound \eqref{5}.  
If Assumption \ref{as:ratio} holds, i.e. $\eta_{\max}=c\eta_{\min}$, then the second term becomes $L c\eta_{\max}\sigma^{2}$, which is absorbed into the $O(\cdot)$ notation used in the theorem statement. \hfill $\square$
\end{enumerate}

\section*{Rate Derivation (With Constants)}
From \eqref{15}, set $\eta_{\max}=c\,\eta_{\min}$ with a fixed ratio $c\ge1$.
Choose $\eta_{\min}= \frac{c_0}{\sqrt{T}}$ where $c_0>0$ satisfies $c_0\le \frac{1}{L\sqrt{c}}$ to respect the stability condition \eqref{3}.
Then
\[
\frac{2\big(f(w_0)-f(w^*)\big)}{\eta_{\min}T}
= \frac{2\big(f(w_0)-f(w^*)\big)}{c_0}\,\frac{1}{\sqrt{T}},
\qquad
\frac{L\,\eta_{\max}^{2}}{\eta_{\min}}\sigma^{2}
= Lc\,c_0\,\frac{1}{\sqrt{T}}\sigma^{2}.
\]
Both terms decay as $O(1/\sqrt{T})$, establishing the $O(1/\sqrt{T})$ convergence rate for the average squared gradient norm.

\section*{Special Cases}
\begin{enumerate}[label=\arabic*.]
\item \textbf{Constant step size.} When $\eta_{\min}=\eta_{\max}=\eta$, the bound \eqref{5} reduces to the classic SGD bound
\[
\frac1T\sum_{t=0}^{T-1}\mathbb{E}\|\nabla f(w_t)\|^{2}
\le \frac{2(f(w_0)-f(w^*))}{\eta T}+L\eta\sigma^{2}.
\]

\item \textbf{Vanishing variance.} If $\sigma=0$ (deterministic gradients), the variance term disappears and we obtain a pure descent rate
\[
\frac1T\sum_{t=0}^{T-1}\mathbb{E}\|\nabla f(w_t)\|^{2}
\le \frac{2(f(w_0)-f(w^*))}{\eta_{\min}T}.
\]

\item \textbf{Large period $M$.} The bound \eqref{5} does not depend explicitly on $M$; however, a larger $M$ yields a slower oscillation, making the effective $\eta_{\min}$ closer to the average step size and often improving practical performance.
\end{enumerate}

\end{document}
% ===== CORRECTION SUMMARY =====
% 1. Step 6 Issue: The original inequality used η_max ≤ 1/L but concluded η_t - (L/2)η_t^2 ≥ η_min/2, which is false without extra work.
%    Fixed by adding Lemma \ref{lem:lb} and using η_t ≤ 1/L to obtain η_t - (L/2)η_t^2 ≥ η_t/2 ≥ η_min/2.
% 2. Step 10 Issue: The transition from (11) to (13) incorrectly simplified Lη_max^2/η_min to Lη_max.
%    Fixed by keeping the exact factor Lη_max^2/η_min in (15) and introducing Assumption \ref{as:ratio} to later bound it by Lcη_max when a constant ratio is assumed.
% 3. Unused bounded‑gradient assumption removed.
% 4. Added explicit stability condition η_max ≤ 1/L to the learning‑rate assumptions.
% 5. Inserted comments and lemmas to keep the proof self‑contained and fully rigorous.